\chapter{La legge dei grandi numeri e il teorema centrale del limite}

Fin qui abbiamo sempre fatto lo stesso tipo di conto: \textbf{nota la legge} di probabilità, calcolare la probabilità
di un evento, una media, una varianza. La statistica fa il cammino inverso: \textbf{osservati dei dati}, risalire alla
legge che li ha prodotti. Con questo capitolo comincia il terzo pacchetto di slide, dedicato alla statistica e alla
misura dell'informazione. Prima di tutto servono due teoremi che spiegano perché fare la media di tanti dati funziona:
la \textbf{legge dei grandi numeri}, che dice che la media dei dati si avvicina alla media vera, e il \textbf{teorema
centrale del limite}, che dice con quale forma (gaussiana) se ne discosta.

% ============================================================
\section{Dalla probabilità alla statistica}
% ============================================================

Supponiamo di avere un \textbf{campione} di dimensione $n$, cioè $n$ numeri $\vett{x} = [x_1, \ldots, x_n]^T \in
\R^n$: le altezze di $n$ persone, $n$ misure di una tensione, i bit ricevuti in un secondo. Il campione è il risultato
di un esperimento aleatorio: se lo si rifacesse (un \emph{ricampionamento}), si otterrebbero numeri diversi.

\dfn{Inferenza statistica}{
  L'\term{inferenza statistica} è il processo che usa l'analisi dei dati per dedurre le proprietà della distribuzione
  di probabilità sottostante, cioè per definire una legge che un qualsiasi campione, scelto a caso, rispetta.
}

L'inferenza si contrappone alla \term{statistica descrittiva}, che si occupa solo delle proprietà dei dati che si hanno
davanti (la loro media, la loro dispersione, una retta che li approssima) senza supporre che provengano da una
popolazione più ampia. Gli obiettivi fondamentali dell'inferenza statistica sono due: i \textbf{test di ipotesi}
(decidere fra due o più spiegazioni dei dati, Capitoli~19 e~20) e la \textbf{stima dei parametri} (trovare il valore
più plausibile di una grandezza non nota, Capitoli~21--23). Alla statistica descrittiva è dedicato il Capitolo~24; gli
ultimi due capitoli misurano l'informazione.

\begin{figure}[H]
\centering
\begin{tikzpicture}[box/.style={draw=mspiA!70!black, fill=mspiA!8, rounded corners=3pt, align=center, font=\small,
  inner sep=6pt}]
  \node[box, text width=5.4cm] (pop) at (0,0) {\textbf{Popolazione}\\[2pt] una legge di probabilità
    $f_X(x; \theta)$, nota solo in parte (per esempio, non si conosce il parametro $\theta$)};
  \node[box, text width=4.8cm, fill=mspiB!8, draw=mspiB!70!black] (cam) at (9.2,0) {\textbf{Campione}\\[2pt] i dati
    osservati $\vett{x} = [x_1, x_2, \ldots, x_n]^T$, uno dei tanti possibili};
  \draw[-{Stealth}, line width=1pt] (pop.north east) to[bend left=16]
    node[above, font=\small, align=center] {esperimento aleatorio\\ (campionamento)} (cam.north west);
  \draw[-{Stealth}, line width=1pt, mspiB!80!black] (cam.south west) to[bend left=16]
    node[below, font=\small, align=center, text=black] {inferenza statistica:\\ test di ipotesi e stima} (pop.south east);
\end{tikzpicture}
\caption{La probabilità va dalla legge ai dati; l'inferenza statistica torna indietro, dai dati alla legge.}
\end{figure}

% ============================================================
\section{La media campione e la legge dei grandi numeri}
% ============================================================

Il primo esempio di inferenza è il più naturale: per conoscere la media di una popolazione si fa la media dei dati.

\dfn{Media campione}{
  Dato un insieme di dati $\vett{x} \in \mathcal{X}^n \subseteq \R^n$, la \term{media campione} è
  \[ \overline{x}_n = \frac{1}{n} \sum_{i=1}^{n} x_i . \]
  Se i dati sono visti come variabili aleatorie, $\vett{X} = [X_1, \ldots, X_n]^T$, la media campione
  $\overline{X}_n = \frac{1}{n} \sum_{i=1}^{n} X_i$ è a sua volta una variabile aleatoria.
}

La convenzione delle lettere è importante: $\overline{x}_n$, con le lettere piccole, è un numero, la media di
\emph{questi} dati (una realizzazione); $\overline{X}_n$, con le lettere grandi, è la media di un campione
\emph{aleatorio}, che cambia da un ricampionamento all'altro. È la media campionaria del Capitolo~5, dove avevamo visto
in una simulazione che, lanciando tante volte un dado, si avvicina a $E[X] = 3{,}5$. Ora lo dimostriamo. Serve
un'ipotesi sul modo in cui sono stati raccolti i dati.

\dfn{Variabili indipendenti e identicamente distribuite}{
  Le variabili $X_1, X_2, \ldots, X_n$ si dicono \term{indipendenti e identicamente distribuite}, in breve
  \term{i.i.d.}, se sono statisticamente indipendenti e hanno tutte la stessa legge (la stessa pmf o la stessa PDF).
}

È il modello delle prove ripetute: ogni dato è prodotto dallo stesso esperimento, rifatto da capo senza che un risultato
influenzi gli altri. In particolare dati i.i.d. hanno tutti la stessa media $\mu = E[X_i]$ e la stessa varianza
$\sigma^2 = E[(X_i - \mu)^2]$.

\teorema[mediacampione]{Media e varianza della media campione}{
  Se $X_1, \ldots, X_n$ sono i.i.d. con media $\mu$ e varianza $\sigma^2$:
  \[ E\left[ \overline{X}_n \right] = E\left[ \frac{X_1 + X_2 + \cdots + X_n}{n} \right] = \mu , \qquad
     \mathrm{VAR}\left[ \overline{X}_n \right] = \mathrm{VAR}\left[ \frac{X_1 + X_2 + \cdots + X_n}{n} \right] =
     \frac{\sigma^2}{n} . \]
}

\begin{dimostrazione}
\passo{La media} Per la linearità della media,
$E[\overline{X}_n] = \frac{1}{n} \left( E[X_1] + \cdots + E[X_n] \right) = \frac{1}{n}\, n\mu = \mu$.
\passo{La varianza} Una costante moltiplicativa esce dalla varianza al quadrato ($\mathrm{VAR}[cY] = c^2\,
\mathrm{VAR}[Y]$, perché la varianza è la media di un quadrato). Le $X_i$ sono indipendenti, quindi incorrelate, e la
varianza della somma è la somma delle varianze (Capitolo~9):
\[ \mathrm{VAR}\left[ \overline{X}_n \right] = \frac{1}{n^2}\, \mathrm{VAR}[X_1 + X_2 + \cdots + X_n] = \frac{1}{n^2}
   \left( \sigma^2 + \sigma^2 + \cdots + \sigma^2 \right) = \frac{n \sigma^2}{n^2} = \frac{\sigma^2}{n} . \]
\end{dimostrazione}

La media campione ha la stessa media dei singoli dati, ma una varianza $n$ volte più piccola: la sua deviazione standard
è $\frac{\sigma}{\sqrt{n}}$. Mediando, gli scarti in su e quelli in giù si compensano in parte. È proprio questo che fa
funzionare il teorema seguente.

\teorema[lgn]{Legge debole dei grandi numeri}{
  Sia $X_1, X_2, \ldots, X_n, \ldots$ una successione di variabili aleatorie i.i.d., con media $E[X_i] = \mu$ e varianza
  $E[(X_i - \mu)^2] = \sigma^2$. Allora, per ogni $\epsilon > 0$,
  \[ \lim_{n \to \infty} P\left( \abs{\frac{X_1 + X_2 + \cdots + X_n}{n} - \mu} > \epsilon \right) = 0 . \]
}

\begin{dimostrazione}[con la disuguaglianza di Chebyshev]
\passo{Chebyshev per la media campione} La disuguaglianza di Chebyshev (Capitolo~7) dice che una variabile si allontana
dalla sua media di almeno $k$ deviazioni standard con probabilità al più $\frac{1}{k^2}$. Per $\overline{X}_n$ la media
è $\mu$ e la deviazione standard è $\frac{\sigma}{\sqrt{n}}$ (teorema precedente); scegliendo $k$ in modo che la
distanza $k\frac{\sigma}{\sqrt{n}}$ sia proprio $\epsilon$, cioè $k = \frac{\epsilon\sqrt{n}}{\sigma}$:
\[ P\left( \abs{\overline{X}_n - \mu} > \epsilon \right) \leq P\left( \abs{\overline{X}_n - \mu} \geq \epsilon \right)
   \leq \frac{1}{k^2} = \frac{\sigma^2}{n\, \epsilon^2} . \]
\passo{Il limite} Il secondo membro tende a 0 per $n \to \infty$ ($\sigma$ ed $\epsilon$ sono fissati), e una
probabilità non può essere negativa: la probabilità, schiacciata fra 0 e una quantità che tende a 0, tende a 0.
\end{dimostrazione}

In parole: comunque si fissi una tolleranza $\epsilon$, piccola quanto si vuole, con abbastanza dati la media campione
cade a meno di $\epsilon$ dalla media vera con probabilità vicina quanto si vuole a 1. La figura mostra sei campioni
diversi di lanci di un dado: ogni media campione segue una strada sua, ma tutte entrano nell'«imbuto» che si stringe
come $\frac{1}{\sqrt{n}}$.

\begin{figure}[H]
\centering
\begin{tikzpicture}
\begin{axis}[width=15cm, height=6.4cm, xmin=0, xmax=400, ymin=1.6, ymax=5.4, xtick={0,50,100,200,300,400},
  ytick={2,2.5,3,3.5,4,4.5,5}, yticklabels={$2$,$2{,}5$,$3$,$3{,}5$,$4$,$4{,}5$,$5$}, xlabel={numero di lanci $n$},
  ylabel={$\overline{x}_n$}, axis lines=left, tick label style={font=\scriptsize}, label style={font=\small},
  grid=major, grid style={mspigray!20}, clip mode=individual]
  \addplot[black, dashed, line width=0.9pt, domain=1:400, samples=200] {3.5 + 2*1.7078/sqrt(x)};
  \addplot[black, dashed, line width=0.9pt, domain=1:400, samples=200] {3.5 - 2*1.7078/sqrt(x)};
  \addplot[black, line width=0.6pt, domain=0:400, samples=2] {3.5};
  \addplot[mspiA, line width=0.8pt] coordinates {(1,2.000) (2,1.500) (3,3.000) (4,3.250) (5,3.200) (6,3.000) (7,2.857) (8,3.000) (9,3.333) (10,3.400) (11,3.273) (12,3.333) (13,3.308) (14,3.357) (15,3.333) (16,3.250) (17,3.235) (18,3.389) (19,3.263) (20,3.250) (25,3.280) (30,3.333) (35,3.400) (40,3.375) (45,3.467) (50,3.540) (55,3.618) (60,3.617) (65,3.646) (70,3.671) (75,3.693) (80,3.675) (85,3.729) (90,3.667) (95,3.621) (100,3.640) (105,3.676) (110,3.664) (115,3.643) (120,3.700) (125,3.688) (130,3.700) (135,3.689) (140,3.721) (145,3.731) (150,3.707) (155,3.735) (160,3.725) (165,3.727) (170,3.688) (175,3.697) (180,3.667) (185,3.643) (190,3.626) (195,3.626) (200,3.640) (205,3.659) (210,3.662) (215,3.637) (220,3.627) (225,3.627) (230,3.639) (235,3.643) (240,3.633) (245,3.629) (250,3.644) (255,3.647) (260,3.642) (265,3.649) (270,3.626) (275,3.636) (280,3.621) (285,3.625) (290,3.641) (295,3.641) (300,3.630) (305,3.649) (310,3.652) (315,3.648) (320,3.666) (325,3.640) (330,3.624) (335,3.633) (340,3.641) (345,3.638) (350,3.626) (355,3.645) (360,3.656) (365,3.663) (370,3.649) (375,3.635) (380,3.616) (385,3.605) (390,3.585) (395,3.585) (400,3.567)};
  \addplot[mspiB, line width=0.8pt] coordinates {(1,5.000) (2,4.000) (3,3.333) (4,2.750) (5,2.600) (6,3.167) (7,3.571) (8,3.375) (9,3.333) (10,3.300) (11,3.545) (12,3.667) (13,3.615) (14,3.714) (15,3.667) (16,3.750) (17,3.824) (18,3.778) (19,3.737) (20,3.700) (25,3.600) (30,3.533) (35,3.543) (40,3.450) (45,3.422) (50,3.460) (55,3.455) (60,3.567) (65,3.538) (70,3.543) (75,3.600) (80,3.587) (85,3.635) (90,3.644) (95,3.632) (100,3.640) (105,3.648) (110,3.682) (115,3.635) (120,3.617) (125,3.592) (130,3.562) (135,3.570) (140,3.521) (145,3.531) (150,3.533) (155,3.548) (160,3.519) (165,3.521) (170,3.518) (175,3.531) (180,3.522) (185,3.551) (190,3.537) (195,3.528) (200,3.550) (205,3.571) (210,3.533) (215,3.512) (220,3.514) (225,3.493) (230,3.491) (235,3.498) (240,3.496) (245,3.498) (250,3.496) (255,3.518) (260,3.500) (265,3.479) (270,3.504) (275,3.498) (280,3.507) (285,3.516) (290,3.510) (295,3.525) (300,3.530) (305,3.511) (310,3.523) (315,3.508) (320,3.500) (325,3.495) (330,3.500) (335,3.501) (340,3.488) (345,3.507) (350,3.497) (355,3.496) (360,3.494) (365,3.479) (370,3.481) (375,3.491) (380,3.487) (385,3.496) (390,3.500) (395,3.519) (400,3.525)};
  \addplot[mspiC, line width=0.8pt] coordinates {(1,2.000) (2,3.500) (3,4.000) (4,3.500) (5,4.000) (6,3.833) (7,3.714) (8,4.000) (9,4.111) (10,3.800) (11,3.818) (12,4.000) (13,3.846) (14,3.714) (15,3.533) (16,3.438) (17,3.353) (18,3.500) (19,3.368) (20,3.350) (25,3.040) (30,3.233) (35,3.114) (40,3.100) (45,3.133) (50,3.260) (55,3.218) (60,3.233) (65,3.138) (70,3.214) (75,3.320) (80,3.275) (85,3.259) (90,3.267) (95,3.232) (100,3.250) (105,3.314) (110,3.318) (115,3.313) (120,3.333) (125,3.416) (130,3.469) (135,3.474) (140,3.457) (145,3.434) (150,3.413) (155,3.413) (160,3.406) (165,3.424) (170,3.400) (175,3.417) (180,3.406) (185,3.454) (190,3.453) (195,3.456) (200,3.450) (205,3.420) (210,3.395) (215,3.363) (220,3.377) (225,3.364) (230,3.378) (235,3.379) (240,3.371) (245,3.376) (250,3.372) (255,3.341) (260,3.342) (265,3.328) (270,3.337) (275,3.316) (280,3.339) (285,3.333) (290,3.331) (295,3.349) (300,3.357) (305,3.377) (310,3.374) (315,3.368) (320,3.394) (325,3.394) (330,3.418) (335,3.415) (340,3.421) (345,3.441) (350,3.437) (355,3.448) (360,3.444) (365,3.466) (370,3.454) (375,3.456) (380,3.466) (385,3.455) (390,3.451) (395,3.451) (400,3.470)};
  \addplot[mspiric, line width=0.8pt] coordinates {(1,2.000) (2,1.500) (3,2.667) (4,2.500) (5,2.200) (6,2.500) (7,2.286) (8,2.125) (9,2.000) (10,2.400) (11,2.364) (12,2.500) (13,2.462) (14,2.714) (15,2.933) (16,3.062) (17,3.176) (18,3.167) (19,3.316) (20,3.400) (25,3.520) (30,3.500) (35,3.457) (40,3.225) (45,3.267) (50,3.260) (55,3.255) (60,3.300) (65,3.369) (70,3.457) (75,3.493) (80,3.487) (85,3.424) (90,3.444) (95,3.442) (100,3.390) (105,3.371) (110,3.445) (115,3.513) (120,3.483) (125,3.456) (130,3.477) (135,3.467) (140,3.471) (145,3.483) (150,3.453) (155,3.465) (160,3.456) (165,3.394) (170,3.365) (175,3.343) (180,3.317) (185,3.324) (190,3.337) (195,3.349) (200,3.340) (205,3.346) (210,3.348) (215,3.353) (220,3.364) (225,3.360) (230,3.374) (235,3.383) (240,3.379) (245,3.412) (250,3.412) (255,3.451) (260,3.469) (265,3.464) (270,3.481) (275,3.480) (280,3.461) (285,3.470) (290,3.472) (295,3.454) (300,3.450) (305,3.443) (310,3.458) (315,3.435) (320,3.431) (325,3.440) (330,3.439) (335,3.433) (340,3.447) (345,3.472) (350,3.480) (355,3.465) (360,3.461) (365,3.479) (370,3.465) (375,3.480) (380,3.482) (385,3.504) (390,3.490) (395,3.501) (400,3.505)};
  \addplot[mspiA!45, line width=0.8pt] coordinates {(1,6.000) (2,6.000) (3,4.667) (4,4.500) (5,4.800) (6,4.333) (7,4.429) (8,4.500) (9,4.222) (10,4.200) (11,4.182) (12,4.000) (13,3.923) (14,3.714) (15,3.800) (16,3.875) (17,3.824) (18,3.944) (19,4.000) (20,4.100) (25,3.960) (30,3.967) (35,3.800) (40,3.800) (45,3.733) (50,3.720) (55,3.691) (60,3.683) (65,3.600) (70,3.643) (75,3.600) (80,3.612) (85,3.635) (90,3.678) (95,3.684) (100,3.620) (105,3.590) (110,3.609) (115,3.600) (120,3.625) (125,3.608) (130,3.608) (135,3.652) (140,3.629) (145,3.662) (150,3.727) (155,3.735) (160,3.712) (165,3.709) (170,3.718) (175,3.686) (180,3.678) (185,3.654) (190,3.674) (195,3.682) (200,3.680) (205,3.668) (210,3.671) (215,3.633) (220,3.618) (225,3.622) (230,3.639) (235,3.651) (240,3.671) (245,3.673) (250,3.704) (255,3.698) (260,3.704) (265,3.694) (270,3.696) (275,3.702) (280,3.686) (285,3.698) (290,3.683) (295,3.692) (300,3.670) (305,3.669) (310,3.661) (315,3.648) (320,3.650) (325,3.634) (330,3.621) (335,3.615) (340,3.624) (345,3.612) (350,3.629) (355,3.634) (360,3.636) (365,3.641) (370,3.630) (375,3.616) (380,3.605) (385,3.605) (390,3.608) (395,3.603) (400,3.612)};
  \addplot[mspiB!50, line width=0.8pt] coordinates {(1,6.000) (2,6.000) (3,5.000) (4,4.500) (5,3.800) (6,4.167) (7,3.857) (8,4.000) (9,4.111) (10,3.900) (11,3.909) (12,3.750) (13,3.769) (14,3.929) (15,3.867) (16,3.688) (17,3.588) (18,3.722) (19,3.579) (20,3.450) (25,3.440) (30,3.300) (35,3.200) (40,3.300) (45,3.378) (50,3.300) (55,3.327) (60,3.483) (65,3.462) (70,3.557) (75,3.520) (80,3.475) (85,3.494) (90,3.478) (95,3.421) (100,3.370) (105,3.343) (110,3.364) (115,3.365) (120,3.383) (125,3.360) (130,3.362) (135,3.378) (140,3.400) (145,3.455) (150,3.467) (155,3.477) (160,3.481) (165,3.497) (170,3.512) (175,3.457) (180,3.417) (185,3.416) (190,3.411) (195,3.415) (200,3.435) (205,3.449) (210,3.424) (215,3.405) (220,3.427) (225,3.427) (230,3.422) (235,3.443) (240,3.475) (245,3.502) (250,3.496) (255,3.490) (260,3.496) (265,3.479) (270,3.467) (275,3.473) (280,3.464) (285,3.449) (290,3.452) (295,3.444) (300,3.447) (305,3.452) (310,3.461) (315,3.454) (320,3.428) (325,3.428) (330,3.442) (335,3.454) (340,3.444) (345,3.458) (350,3.466) (355,3.465) (360,3.475) (365,3.479) (370,3.508) (375,3.493) (380,3.513) (385,3.506) (390,3.508) (395,3.527) (400,3.540)};
  \node[font=\scriptsize, anchor=west] at (axis cs:404,3.5) {$\mu = 3{,}5$};
  \node[font=\scriptsize, anchor=south west, inner sep=2pt] at (axis cs:40,4.04) {$\mu + 2\frac{\sigma}{\sqrt{n}}$};
  \node[font=\scriptsize, anchor=north west, inner sep=2pt] at (axis cs:40,2.96) {$\mu - 2\frac{\sigma}{\sqrt{n}}$};
\end{axis}
\end{tikzpicture}
\caption{Sei campioni di lanci di un dado onesto (simulazione) e le loro medie campione $\overline{x}_n$ al crescere di
$n$. Le curve tratteggiate distano due deviazioni standard della media campione, $2\frac{\sigma}{\sqrt{n}}$ con
$\sigma = \sqrt{35/12} \approx 1{,}71$, dalla media vera: l'imbuto si stringe e le medie vi restano dentro quasi sempre.}
\end{figure}

\begin{esempio}{Quanti lanci servono?}
Si vuole stimare la probabilità di «testa» di una moneta onesta con la frequenza relativa, e si vuole che l'errore
superi $0{,}05$ con probabilità al più del 5\%. Con $X_i = 1$ se esce testa e $0$ altrimenti, $\mu = \frac12$ e
$\sigma^2 = E[X^2] - \mu^2 = \frac12 - \frac14 = \frac14$. Per la dimostrazione della legge dei grandi numeri
\[ P\left( \abs{\overline{X}_n - \tfrac12} > 0{,}05 \right) \leq \frac{\sigma^2}{n\, \epsilon^2} = \frac{1/4}{n \cdot
   0{,}0025} = \frac{100}{n} , \]
che è al più $0{,}05$ se $n \geq 2000$. Duemila lanci bastano di sicuro, qualunque cosa succeda. Vedremo alla fine del
capitolo che, grazie al teorema centrale del limite, ne bastano in realtà circa 400: Chebyshev vale per qualsiasi legge,
e per questo è prudente.
\end{esempio}

\subsection{Le frequenze tendono alle probabilità}

La legge dei grandi numeri giustifica l'interpretazione frequentista della probabilità da cui eravamo partiti nel
Capitolo~3.

\teorema[frequenze]{La frequenza relativa tende alla probabilità}{
  Si ripetano $n$ copie indipendenti di un esperimento, in ciascuna delle quali può verificarsi un evento $E$. La
  frazione $f_n(E)$ delle $n$ prove in cui $E$ si è verificato tende alla probabilità $P(E)$, nel senso della legge
  debole dei grandi numeri: per ogni $\epsilon > 0$, $\lim_{n \to \infty} P\left( \abs{f_n(E) - P(E)} > \epsilon
  \right) = 0$.
}

\begin{dimostrazione}[con le variabili indicatrici]
\passo{Una variabile che conta} Si pone
\[ X_i = \begin{cases} 1 & \text{se } E \text{ si realizza nell'esperimento } i\text{-esimo} , \\ 0 & \text{se } E
   \text{ non si realizza nell'esperimento } i\text{-esimo} . \end{cases} \]
La somma $X_1 + X_2 + \cdots + X_n$ conta le prove, fra le prime $n$, in cui si è verificato $E$; divisa per $n$ è
proprio la frequenza relativa: $f_n(E) = \overline{X}_n$.
\passo{Media e varianza} Le $X_i$ sono i.i.d. (le prove sono copie indipendenti dello stesso esperimento), con
$E[X_i] = 1 \cdot P(E) + 0 \cdot P(\compl{E}) = P(E)$ e varianza finita $P(E)\,(1 - P(E))$ (come per la Bernoulli del
Capitolo~7).
\passo{Si applica la legge dei grandi numeri} con $\mu = P(E)$.
\end{dimostrazione}

% ============================================================
\section{Modi diversi di convergere}
% ============================================================

Per una successione di numeri «tendere a un limite» ha un solo significato (Capitolo~3). Per una successione di
\emph{variabili aleatorie}, come le medie campione $\overline{X}_1, \overline{X}_2, \ldots$, ce ne sono diversi.
\begin{itemize}
  \item \textbf{Convergenza in probabilità} (o \emph{convergenza debole}): per ogni $\epsilon > 0$,
        $\lim_{n \to \infty} P\left( \abs{\overline{X}_n - E[X]} > \epsilon \right) = 0$. È quella della legge debole
        dei grandi numeri. Si noti che è il limite di una successione di \textbf{numeri}, le probabilità: una
        successione deterministica.
  \item \textbf{Convergenza quasi certa} (o \emph{forte}, o \emph{con probabilità uno}): la probabilità che la
        successione aleatoria delle medie campione converga a $E[X]$ vale 1,
        $P\left( \lim_{n \to \infty} \overline{X}_n = E[X] \right) = 1$. Al limite, la probabilità di allontanarsi
        dalla media statistica è zero. È la \term{legge forte dei grandi numeri}, che vale anch'essa per variabili
        i.i.d. (non la dimostriamo). Qui l'ordine è rovesciato: è la probabilità di un limite, cioè la probabilità
        che la singola realizzazione della successione, la singola curva della figura, converga.
  \item \textbf{Convergenza in media quadratica}: $\lim_{n \to \infty} E\left[ \left( \overline{X}_n - E[X]
        \right)^2 \right] = 0$. Per la media campione è immediata: quella media è proprio la varianza di
        $\overline{X}_n$, cioè $\frac{\sigma^2}{n} \to 0$.
\end{itemize}

\begin{center}
\begin{tikzpicture}[box/.style={draw=mspigray, rounded corners=3pt, align=center, font=\small, inner sep=5pt,
  text width=4.4cm, fill=mspilight}]
  \node[box] (qc) at (0,1.15) {convergenza quasi certa\\ (legge forte)};
  \node[box] (mq) at (0,-1.15) {convergenza in media quadratica};
  \node[box, fill=evid!60] (pr) at (7.5,0) {convergenza in probabilità\\ (legge debole)};
  \draw[-{Stealth}, line width=1pt] (qc.east) -- node[above, font=\scriptsize, sloped] {implica} (pr.north west);
  \draw[-{Stealth}, line width=1pt] (mq.east) -- node[below, font=\scriptsize, sloped] {implica} (pr.south west);
\end{tikzpicture}
\end{center}

Le due convergenze di sinistra sono più forti e implicano quella di destra. Per la media quadratica lo si vede con la
disuguaglianza di Markov, come nella legge dei grandi numeri: se $Y_n$ è una successione di variabili e $c$ un numero,
\[ P\left( \abs{Y_n - c} > \epsilon \right) = P\left( (Y_n - c)^2 > \epsilon^2 \right) \leq \frac{E\left[ (Y_n -
   c)^2 \right]}{\epsilon^2} , \]
e se il numeratore tende a zero anche la probabilità tende a zero.

\begin{esempio}{Le frequenze di un alfabeto finito}
I dati $\vett{x} \in \mathcal{X}^n$ abbiano un alfabeto discreto e finito, $\mathcal{X} = \{a_1, \ldots, a_M\}$, e le
$X_k$ siano i.i.d. con pmf $p_X$. Sia $N_i$ il numero di volte in cui $X_k = a_i$ nel campione. Ogni dato vale $a_i$ con
probabilità $p_X(a_i)$, indipendentemente dagli altri, quindi $N_i$ conta i successi in $n$ prove ed è binomiale
(Capitolo~5):
\[ P(N_i = k) = \binom{n}{k}\, p_X(a_i)^k\, \left[ 1 - p_X(a_i) \right]^{n - k} . \]
La frequenza relativa $\frac{N_i}{n}$ è anch'essa una variabile aleatoria. Con la media $n\,p_X(a_i)$ e la varianza
$n\,p_X(a_i)(1 - p_X(a_i))$ della binomiale (Capitolo~7):
\[ E\left[ \frac{N_i}{n} \right] = p_X(a_i) , \qquad \mathrm{VAR}\left[ \frac{N_i}{n} \right] = \frac{p_X(a_i) \left(
   1 - p_X(a_i) \right)}{n} . \]
Siccome la media di $\frac{N_i}{n}$ è proprio $p_X(a_i)$, la varianza è la media dello scarto quadratico dalla
probabilità, e quindi
\[ \lim_{n \to \infty} E\left[ \left( \frac{N_i}{n} - p_X(a_i) \right)^2 \right] = \lim_{n \to \infty} \frac{p_X(a_i)
   \left( 1 - p_X(a_i) \right)}{n} = 0 : \]
la frequenza relativa converge alla probabilità in media quadratica (e quindi anche in probabilità).
\end{esempio}

% ============================================================
\section{Il teorema centrale del limite}
% ============================================================

La legge dei grandi numeri dice \emph{dove} va a finire la media campione. Il teorema centrale del limite dice qualcosa
di più sorprendente: \emph{con quale forma} si distribuiscono le somme di tante variabili.

Si consideri la somma $Z_n = \sum_{i=1}^{n} X_i$ di $n$ variabili i.i.d., con densità $f_{X_i}(x) = f_X(x)$, media
$\mu_X$ e varianza $\sigma_X^2$. Per la linearità della media e per l'indipendenza
\[ \mu_{Z_n} = n\, \mu_X , \qquad \sigma_{Z_n}^2 = n\, \sigma_X^2 : \]
la somma ha un centro che si sposta e una larghezza che cresce. Per confrontare somme con $n$ diversi conviene
standardizzarle (Capitolo~15): la variabile
\[ S_n = \frac{Z_n - \mu_{Z_n}}{\sigma_{Z_n}} = \frac{Z_n - n\, \mu_X}{\sqrt{n}\, \sigma_X} \]
ha sempre, indipendentemente da $n$, valore atteso nullo e varianza unitaria: è una variabile standard.

\teorema[tcl]{Teorema centrale del limite}{
  Sotto ipotesi poco restrittive sulla densità $f_X(x)$ delle $X_i$ (per esempio, varianza finita), al tendere di $n$
  all'infinito la CDF di $S_n$ tende alla CDF di una gaussiana standard:
  \[ \lim_{n \to \infty} F_{S_n}(x) = \lim_{n \to \infty} P(S_n \leq x) = \Phi(x) = \int_{-\infty}^{x} \frac{1}{\sqrt{2\pi}}
     e^{-\frac{t^2}{2}} \dd t = 1 - Q(x) . \]
  In sostanza, la somma di un gran numero di variabili aleatorie statisticamente indipendenti segue con buona
  approssimazione una legge gaussiana, \textbf{a prescindere dalla particolare distribuzione di ciascuna di esse}.
}

Le slide osservano che l'ipotesi di equidistribuzione permette una dimostrazione particolarmente semplice, e che il
teorema vale anche sotto ipotesi meno restrittive (variabili non identiche), con una dimostrazione più complessa. Qui ne
diamo l'idea, con uno strumento nuovo; due fatti tecnici vanno accettati senza dimostrazione.

\begin{dimostrazione}[l'idea, con la funzione generatrice dei momenti]
\passo{Uno strumento: la funzione generatrice dei momenti} Per una variabile $Y$ si definisce $M_Y(t) = E\left[ e^{tY}
\right]$, una funzione del numero $t$. Si accettano due fatti: se due variabili hanno la stessa $M$ (per $t$ vicino a
0) hanno la stessa legge; e se $M_{Y_n}(t) \to M_Y(t)$, allora le CDF di $Y_n$ tendono a quella di $Y$.
\passo{La $M$ della gaussiana standard} Completando il quadrato, $tx - \frac{x^2}{2} = -\frac{(x - t)^2}{2} +
\frac{t^2}{2}$, quindi
\[ M(t) = \int_{-\infty}^{\infty} e^{tx}\, \frac{1}{\sqrt{2\pi}} e^{-\frac{x^2}{2}} \dd x = e^{\frac{t^2}{2}}
   \int_{-\infty}^{\infty} \frac{1}{\sqrt{2\pi}} e^{-\frac{(x - t)^2}{2}} \dd x = e^{\frac{t^2}{2}} , \]
perché l'ultimo integrale è l'area di una gaussiana di media $t$, cioè 1.
\passo{La $M$ di $S_n$} Siano $Y_i = \frac{X_i - \mu_X}{\sigma_X}$ i dati standardizzati: sono i.i.d., con media 0 e
varianza 1, e $S_n = \frac{1}{\sqrt{n}} \sum_i Y_i$. L'esponenziale di una somma è il prodotto degli esponenziali, e
la media di un prodotto di variabili indipendenti è il prodotto delle medie (Capitolo~9):
\[ M_{S_n}(t) = E\left[ e^{\frac{t}{\sqrt{n}} Y_1} \cdots e^{\frac{t}{\sqrt{n}} Y_n} \right] = E\left[
   e^{\frac{t}{\sqrt{n}} Y_1} \right] \cdots E\left[ e^{\frac{t}{\sqrt{n}} Y_n} \right] = \left[ M_Y\left(
   \frac{t}{\sqrt{n}} \right) \right]^n . \]
\passo{Per $n$ grande l'argomento è piccolo} Con la serie esponenziale $e^u = 1 + u + \frac{u^2}{2} + \cdots$
(Capitolo~5) e la linearità della media, ricordando $E[Y] = 0$ ed $E[Y^2] = 1$:
\[ M_Y\left( \frac{t}{\sqrt{n}} \right) = 1 + \frac{t}{\sqrt{n}}\, E[Y] + \frac{t^2}{2n}\, E[Y^2] + \cdots = 1 +
   \frac{t^2}{2n} + (\text{termini che vanno a zero più in fretta di } \tfrac{1}{n}) . \]
Tutte le informazioni sulla legge di $Y$ oltre a media e varianza finiscono nei termini trascurabili: per questo il
risultato non dipende dalla legge di partenza.
\passo{Il limite} Per il limite $\left( 1 + \frac{x}{n} \right)^n \to e^x$ (Capitolo~5, con $x = \frac{t^2}{2}$),
\[ M_{S_n}(t) = \left[ 1 + \frac{t^2}{2n} + \cdots \right]^n \;\longrightarrow\; e^{\frac{t^2}{2}} , \]
che è la $M$ della gaussiana standard: per il secondo fatto del passo 1, $F_{S_n}(x) \to \Phi(x)$.
\end{dimostrazione}

\info{Perché solo un'idea}{
  La funzione generatrice dei momenti non esiste per tutte le variabili (per la Cauchy, per esempio, la media di
  $e^{tY}$ è infinita). La dimostrazione completa usa la \emph{funzione caratteristica} $E\left[ e^{jtY} \right]$, con
  l'unità immaginaria $j$, che esiste sempre: i passaggi sono gli stessi. I due fatti accettati nel passo 1 sono teoremi
  di analisi che richiedono strumenti oltre questo corso.
}

Le slide aggiungono un'osservazione pratica: sommando $N$ variabili uniformi i.i.d., già per $N = 15$ si è vicini a una
densità normale, ma la convergenza è rapida attorno alla media e lenta sulle code. La figura lo mostra: a sinistra le
densità di $S_N$ per $N = 1, 2, 3$ (la somma di due uniformi ha densità triangolare, di tre è fatta di archi di
parabola), a destra il confronto sulle code per $N = 15$, in scala logaritmica.

\begin{figure}[H]
\centering
\begin{tikzpicture}
\begin{groupplot}[group style={group size=2 by 1, horizontal sep=1.6cm}, height=5.8cm,
  tick label style={font=\scriptsize}, label style={font=\small}]
\nextgroupplot[width=8.4cm, xmin=-3.8, xmax=3.8, ymin=0, ymax=0.47, xtick={-3,-2,-1,0,1,2,3}, ytick={0.1,0.2,0.3,0.4},
  yticklabels={$0{,}1$,$0{,}2$,$0{,}3$,$0{,}4$}, xlabel=$s$, axis lines=left, grid=major, grid style={mspigray!20},
  legend style={font=\scriptsize, at={(0.02,0.98)}, anchor=north west}, title={somme di $N$ uniformi, standardizzate}]
  \addplot[mspiC, line width=1pt] coordinates {(-3.8,0) (-1.732,0) (-1.732,0.2887) (1.732,0.2887) (1.732,0) (3.8,0)};
  \addplot[mspiB, line width=1pt] coordinates {(-4.00,0.00000) (-3.95,0.00000) (-3.90,0.00000) (-3.85,0.00000) (-3.80,0.00000) (-3.75,0.00000) (-3.70,0.00000) (-3.65,0.00000) (-3.60,0.00000) (-3.55,0.00000) (-3.50,0.00000) (-3.45,0.00000) (-3.40,0.00000) (-3.35,0.00000) (-3.30,0.00000) (-3.25,0.00000) (-3.20,0.00000) (-3.15,0.00000) (-3.10,0.00000) (-3.05,0.00000) (-3.00,0.00000) (-2.95,0.00000) (-2.90,0.00000) (-2.85,0.00000) (-2.80,0.00000) (-2.75,0.00000) (-2.70,0.00000) (-2.65,0.00000) (-2.60,0.00000) (-2.55,0.00000) (-2.50,0.00000) (-2.45,0.00000) (-2.40,0.00825) (-2.35,0.01658) (-2.30,0.02491) (-2.25,0.03325) (-2.20,0.04158) (-2.15,0.04991) (-2.10,0.05825) (-2.05,0.06658) (-2.00,0.07491) (-1.95,0.08325) (-1.90,0.09158) (-1.85,0.09991) (-1.80,0.10825) (-1.75,0.11658) (-1.70,0.12491) (-1.65,0.13325) (-1.60,0.14158) (-1.55,0.14991) (-1.50,0.15825) (-1.45,0.16658) (-1.40,0.17491) (-1.35,0.18325) (-1.30,0.19158) (-1.25,0.19991) (-1.20,0.20825) (-1.15,0.21658) (-1.10,0.22491) (-1.05,0.23325) (-1.00,0.24158) (-0.95,0.24991) (-0.90,0.25825) (-0.85,0.26658) (-0.80,0.27491) (-0.75,0.28325) (-0.70,0.29158) (-0.65,0.29991) (-0.60,0.30825) (-0.55,0.31658) (-0.50,0.32491) (-0.45,0.33325) (-0.40,0.34158) (-0.35,0.34991) (-0.30,0.35825) (-0.25,0.36658) (-0.20,0.37491) (-0.15,0.38325) (-0.10,0.39158) (-0.05,0.39991) (0.00,0.40825) (0.05,0.39991) (0.10,0.39158) (0.15,0.38325) (0.20,0.37491) (0.25,0.36658) (0.30,0.35825) (0.35,0.34991) (0.40,0.34158) (0.45,0.33325) (0.50,0.32491) (0.55,0.31658) (0.60,0.30825) (0.65,0.29991) (0.70,0.29158) (0.75,0.28325) (0.80,0.27491) (0.85,0.26658) (0.90,0.25825) (0.95,0.24991) (1.00,0.24158) (1.05,0.23325) (1.10,0.22491) (1.15,0.21658) (1.20,0.20825) (1.25,0.19991) (1.30,0.19158) (1.35,0.18325) (1.40,0.17491) (1.45,0.16658) (1.50,0.15825) (1.55,0.14991) (1.60,0.14158) (1.65,0.13325) (1.70,0.12491) (1.75,0.11658) (1.80,0.10825) (1.85,0.09991) (1.90,0.09158) (1.95,0.08325) (2.00,0.07491) (2.05,0.06658) (2.10,0.05825) (2.15,0.04991) (2.20,0.04158) (2.25,0.03325) (2.30,0.02491) (2.35,0.01658) (2.40,0.00825) (2.45,0.00000) (2.50,0.00000) (2.55,0.00000) (2.60,0.00000) (2.65,0.00000) (2.70,0.00000) (2.75,0.00000) (2.80,0.00000) (2.85,0.00000) (2.90,0.00000) (2.95,0.00000) (3.00,0.00000) (3.05,0.00000) (3.10,0.00000) (3.15,0.00000) (3.20,0.00000) (3.25,0.00000) (3.30,0.00000) (3.35,0.00000) (3.40,0.00000) (3.45,0.00000) (3.50,0.00000) (3.55,0.00000) (3.60,0.00000) (3.65,0.00000) (3.70,0.00000) (3.75,0.00000) (3.80,0.00000) (3.85,0.00000) (3.90,0.00000) (3.95,0.00000) (4.00,0.00000)};
  \addplot[mspiA, line width=1pt] coordinates {(-4.00,0.00000) (-3.95,0.00000) (-3.90,0.00000) (-3.85,0.00000) (-3.80,0.00000) (-3.75,0.00000) (-3.70,0.00000) (-3.65,0.00000) (-3.60,0.00000) (-3.55,0.00000) (-3.50,0.00000) (-3.45,0.00000) (-3.40,0.00000) (-3.35,0.00000) (-3.30,0.00000) (-3.25,0.00000) (-3.20,0.00000) (-3.15,0.00000) (-3.10,0.00000) (-3.05,0.00000) (-3.00,0.00000) (-2.95,0.00016) (-2.90,0.00063) (-2.85,0.00141) (-2.80,0.00250) (-2.75,0.00391) (-2.70,0.00562) (-2.65,0.00766) (-2.60,0.01000) (-2.55,0.01266) (-2.50,0.01562) (-2.45,0.01891) (-2.40,0.02250) (-2.35,0.02641) (-2.30,0.03063) (-2.25,0.03516) (-2.20,0.04000) (-2.15,0.04516) (-2.10,0.05062) (-2.05,0.05641) (-2.00,0.06250) (-1.95,0.06891) (-1.90,0.07563) (-1.85,0.08266) (-1.80,0.09000) (-1.75,0.09766) (-1.70,0.10563) (-1.65,0.11391) (-1.60,0.12250) (-1.55,0.13141) (-1.50,0.14062) (-1.45,0.15016) (-1.40,0.16000) (-1.35,0.17016) (-1.30,0.18062) (-1.25,0.19141) (-1.20,0.20250) (-1.15,0.21391) (-1.10,0.22562) (-1.05,0.23766) (-1.00,0.25000) (-0.95,0.26219) (-0.90,0.27375) (-0.85,0.28469) (-0.80,0.29500) (-0.75,0.30469) (-0.70,0.31375) (-0.65,0.32219) (-0.60,0.33000) (-0.55,0.33719) (-0.50,0.34375) (-0.45,0.34969) (-0.40,0.35500) (-0.35,0.35969) (-0.30,0.36375) (-0.25,0.36719) (-0.20,0.37000) (-0.15,0.37219) (-0.10,0.37375) (-0.05,0.37469) (0.00,0.37500) (0.05,0.37469) (0.10,0.37375) (0.15,0.37219) (0.20,0.37000) (0.25,0.36719) (0.30,0.36375) (0.35,0.35969) (0.40,0.35500) (0.45,0.34969) (0.50,0.34375) (0.55,0.33719) (0.60,0.33000) (0.65,0.32219) (0.70,0.31375) (0.75,0.30469) (0.80,0.29500) (0.85,0.28469) (0.90,0.27375) (0.95,0.26219) (1.00,0.25000) (1.05,0.23766) (1.10,0.22563) (1.15,0.21391) (1.20,0.20250) (1.25,0.19141) (1.30,0.18063) (1.35,0.17016) (1.40,0.16000) (1.45,0.15016) (1.50,0.14062) (1.55,0.13141) (1.60,0.12250) (1.65,0.11391) (1.70,0.10562) (1.75,0.09766) (1.80,0.09000) (1.85,0.08266) (1.90,0.07563) (1.95,0.06891) (2.00,0.06250) (2.05,0.05641) (2.10,0.05063) (2.15,0.04516) (2.20,0.04000) (2.25,0.03516) (2.30,0.03063) (2.35,0.02641) (2.40,0.02250) (2.45,0.01891) (2.50,0.01562) (2.55,0.01266) (2.60,0.01000) (2.65,0.00766) (2.70,0.00562) (2.75,0.00391) (2.80,0.00250) (2.85,0.00141) (2.90,0.00062) (2.95,0.00016) (3.00,0.00000) (3.05,0.00000) (3.10,0.00000) (3.15,0.00000) (3.20,0.00000) (3.25,0.00000) (3.30,0.00000) (3.35,0.00000) (3.40,0.00000) (3.45,0.00000) (3.50,0.00000) (3.55,0.00000) (3.60,0.00000) (3.65,0.00000) (3.70,0.00000) (3.75,0.00000) (3.80,0.00000) (3.85,0.00000) (3.90,0.00000) (3.95,0.00000) (4.00,0.00000)};
  \addplot[black, dashed, line width=1pt, domain=-3.8:3.8, samples=121] {exp(-x^2/2)/sqrt(2*pi)};
  \legend{{$N = 1$}, {$N = 2$}, {$N = 3$}, {gaussiana}}
\nextgroupplot[width=7.4cm, xmin=0, xmax=6.5, ymin=1e-9, ymax=1, ymode=log, xtick={0,1,2,3,4,5,6}, xlabel=$s$,
  axis lines=left, grid=major, grid style={mspigray!20}, legend style={font=\scriptsize, at={(0.02,0.03)},
  anchor=south west}, title={le code, per $N = 15$}]
  \addplot[mspiA, line width=1pt] coordinates {(0.0,3.9493e-01) (0.1,3.9304e-01) (0.2,3.8742e-01) (0.3,3.7822e-01) (0.4,3.6570e-01) (0.5,3.5019e-01) (0.6,3.3210e-01) (0.7,3.1189e-01) (0.8,2.9007e-01) (0.9,2.6712e-01) (1.0,2.4357e-01) (1.1,2.1988e-01) (1.2,1.9652e-01) (1.3,1.7386e-01) (1.4,1.5224e-01) (1.5,1.3193e-01) (1.6,1.1314e-01) (1.7,9.6005e-02) (1.8,8.0589e-02) (1.9,6.6913e-02) (2.0,5.4945e-02) (2.1,4.4610e-02) (2.2,3.5806e-02) (2.3,2.8405e-02) (2.4,2.2265e-02) (2.5,1.7241e-02) (2.6,1.3185e-02) (2.7,9.9554e-03) (2.8,7.4187e-03) (2.9,5.4544e-03) (3.0,3.9550e-03) (3.1,2.8271e-03) (3.2,1.9913e-03) (3.3,1.3813e-03) (3.4,9.4321e-04) (3.5,6.3357e-04) (3.6,4.1837e-04) (3.7,2.7139e-04) (3.8,1.7279e-04) (3.9,1.0789e-04) (4.0,6.5993e-05) (4.1,3.9500e-05) (4.2,2.3105e-05) (4.3,1.3189e-05) (4.4,7.3348e-06) (4.5,3.9665e-06) (4.6,2.0812e-06) (4.7,1.0569e-06) (4.8,5.1798e-07) (4.9,2.4391e-07) (5.0,1.1008e-07) (5.1,4.7322e-08) (5.2,1.9251e-08) (5.3,7.2065e-09) (5.4,2.5698e-09) (5.5,1.0697e-09) (5.6,1.2430e-10) (5.7,2.6789e-10) (5.8,6.7505e-11) (5.9,1.0000e-12) (6.0,2.4907e-10) (6.1,4.9800e-10) (6.2,2.8847e-10) (6.3,4.7899e-10) (6.4,1.0000e-12) (6.5,6.0975e-11)};
  \addplot[black, dashed, line width=1pt, domain=0:6.5, samples=80] {exp(-x^2/2)/sqrt(2*pi)};
  \legend{{$N = 15$}, {gaussiana}}
\end{groupplot}
\end{tikzpicture}
\caption{Il teorema centrale del limite in azione. A sinistra: la densità della somma standardizzata di $N$ uniformi
diventa una campana già con $N = 3$. A destra, in scala logaritmica: con $N = 15$ le due curve coincidono vicino al
centro, ma a $s = 5$ differiscono di un fattore 13, e oltre $s = \sqrt{45} \approx 6{,}7$ la somma di 15 uniformi non
arriva proprio, mentre la gaussiana sì.}
\end{figure}

\begin{esempio}{La binomiale diventa gaussiana}
Il numero di teste in $n$ lanci di una moneta è una somma di $n$ variabili di Bernoulli i.i.d.: per il teorema centrale
del limite, per $n$ grande è circa gaussiano, con la stessa media $np$ e la stessa varianza $np(1 - p)$ della binomiale.
Per $n = 50$ e $p = \frac12$, la pmf binomiale e la densità $\mathcal{N}(25;\ 12{,}5)$ sono quasi indistinguibili.
\begin{center}
\begin{tikzpicture}
\begin{axis}[quaderno, axis lines=left, width=12.5cm, height=5cm, xmin=9, xmax=41, ymin=0, ymax=0.125,
  xtick={10,15,20,25,30,35,40}, ytick={0.05,0.1}, yticklabels={$0{,}05$,$0{,}1$}, xlabel=$k$,
  legend style={font=\scriptsize, at={(0.98,0.98)}, anchor=north east}]
  \addplot[bastoncini, mspiA, mark size=1.3pt, line width=1.2pt] coordinates {(10,0.00001) (11,0.00003) (12,0.00011) (13,0.00032) (14,0.00083) (15,0.00200) (16,0.00437) (17,0.00875) (18,0.01603) (19,0.02701) (20,0.04186) (21,0.05980) (22,0.07883) (23,0.09596) (24,0.10796) (25,0.11228) (26,0.10796) (27,0.09596) (28,0.07883) (29,0.05980) (30,0.04186) (31,0.02701) (32,0.01603) (33,0.00875) (34,0.00437) (35,0.00200) (36,0.00083) (37,0.00032) (38,0.00011) (39,0.00003) (40,0.00001)};
  \addplot[mspiB, line width=1pt, domain=9:41, samples=120] {exp(-(x-25)^2/25)/sqrt(25*pi)};
  \legend{{$P(\text{$k$ teste in 50 lanci})$}, {densità $\mathcal{N}(25;\ 12{,}5)$}}
\end{axis}
\end{tikzpicture}
\end{center}
Torniamo all'esempio dei lanci di una moneta. Per il teorema centrale del limite $\overline{X}_n$ è circa gaussiana,
con media $\frac12$ e deviazione standard $\frac{\sigma}{\sqrt{n}} = \frac{1}{2\sqrt{n}}$; con la regola delle code
del Capitolo~15 (due code uguali):
\[ P\left( \abs{\overline{X}_n - \tfrac12} > 0{,}05 \right) \approx 2\, Q\left( \frac{0{,}05}{1/(2\sqrt{n})} \right) =
   2\, Q\left( 0{,}1\, \sqrt{n} \right) . \]
Questo vale $0{,}05$ quando $Q(0{,}1\sqrt{n}) = 0{,}025$, cioè $0{,}1\sqrt{n} \approx 1{,}96$ (dalle tavole della
funzione $Q$): $n \approx 385$, contro i 2000 garantiti da Chebyshev.
\end{esempio}

\warning{Quando il teorema non vale}{
  Servono la media e la varianza finite. Per la variabile di Cauchy (Capitoli~11 e~13), che non ha né l'una né l'altra,
  la media di $n$ dati i.i.d. è di nuovo una Cauchy, larga come un dato singolo: né la legge dei grandi numeri né il
  teorema centrale del limite funzionano, come si vedeva nella figura delle medie progressive del Capitolo~11.
}

\info{Perché la gaussiana è ovunque}{
  Il teorema centrale del limite spiega l'onnipresenza della gaussiana annunciata nel Capitolo~15. Il rumore termico di
  un resistore è la somma degli effetti di moltissimi elettroni; un errore di misura è la somma di tante piccole cause
  indipendenti; l'altezza di una persona dipende da molti geni e da molti fattori ambientali. Quale che sia la legge
  dei singoli contributi, la somma è gaussiana.
}

\gbox{In sintesi}{mspiA!80!black}{
\begin{itemize}
  \item L'\textbf{inferenza statistica} risale dai dati alla legge che li ha prodotti; gli obiettivi sono i test di
        ipotesi e la stima. La \textbf{statistica descrittiva} guarda solo i dati.
  \item Per dati i.i.d., la \textbf{media campione} $\overline{X}_n$ ha media $\mu$ e varianza $\frac{\sigma^2}{n}$.
  \item \textbf{Legge debole dei grandi numeri}: $P(\abs{\overline{X}_n - \mu} > \epsilon) \leq \frac{\sigma^2}{n
        \epsilon^2} \to 0$ (Chebyshev). Le frequenze relative tendono alle probabilità.
  \item Convergenza in probabilità (legge debole), quasi certa (legge forte), in media quadratica; le ultime due
        implicano la prima.
  \item \textbf{Teorema centrale del limite}: la somma standardizzata $S_n = \frac{Z_n - n\mu_X}{\sqrt{n}\, \sigma_X}$ di
        $n$ variabili i.i.d. ha CDF che tende a $\Phi(x) = 1 - Q(x)$, qualunque sia la legge di partenza (con varianza
        finita). La convergenza è rapida al centro e lenta sulle code.
\end{itemize}
}
